% VI29 detailed challenge solutions -- VI/24+ proof-depth standard.

% VI29-DETAILED-CHALLENGE-01
\subsection*{Challenge 1 solution}
\begin{solution}
For a Noetherian affine scheme \(X=\Spec A\), define the height strata
\[
H_r(X)=\{\mathfrak p\in\Spec A:\operatorname{ht}(\mathfrak p)=r\}.
\]
A specialization
\[
\mathfrak p\rightsquigarrow\mathfrak q
\]
means \(\mathfrak p\subseteq\mathfrak q\).  Every strict specialization increases height by at
least one, because any chain ending at \(\mathfrak p\) can be extended by the strict inclusion
\(\mathfrak p\subsetneq\mathfrak q\).  Thus the height function measures how far a point can lie
below a terminal point in the specialization order.

For
\[
X=\Spec k[x,y],
\]
the polynomial-ring dimension theorem gives \(\dim X=2\).  There are three possible heights:
\[
H_0,\qquad H_1,\qquad H_2.
\]

Because \(k[x,y]\) is a domain, its unique minimal prime is \((0)\), so
\[
H_0=\{(0)\}.
\]
This is the generic point.

Height-one primes are the nonzero principal prime ideals
\[
(f),
\]
where \(f\in k[x,y]\) is irreducible.  Indeed \(k[x,y]\) is a UFD, so every height-one prime is
principal, while the principal ideal theorem gives height at most \(1\) and
\((0)\subsetneq(f)\) gives height at least \(1\).  Geometrically these are the irreducible curves
in the affine plane.

Maximal ideals have height \(2\).  Over an algebraically closed field they are
\[
(x-a,y-b),
\]
and the chain
\[
(0)\subsetneq(x-a)\subsetneq(x-a,y-b)
\]
shows height \(2\).  Over an arbitrary field the same height statement follows from the standard
dimension theory of the two-variable polynomial ring.  These are the closed points.

Thus a maximal specialization chain has the form
\[
\text{generic point}
\rightsquigarrow
\text{irreducible-curve point}
\rightsquigarrow
\text{closed point},
\]
corresponding algebraically to
\[
(0)\subsetneq(f)\subsetneq\mathfrak m.
\]
The three height strata therefore record the three levels of the specialization order in the
affine plane.
\end{solution}

% VI29-DETAILED-CHALLENGE-02
\subsection*{Challenge 2 solution}
\begin{solution}
Take
\[
A=k[x,y]\times k[t].
\]
This is a finite-type \(k\)-algebra: it is generated, for example, by
\[
(x,0),\ (y,0),\ (0,t),\ (1,0)
\]
together with the diagonal copy of \(k\).

Its spectrum is the disjoint union
\[
\Spec A
=
\Spec k[x,y]\sqcup\Spec k[t].
\]
Hence there are exactly two irreducible components.  The first has dimension
\[
\dim k[x,y]=2,
\]
and the second has dimension
\[
\dim k[t]=1.
\]
A prime chain cannot cross between the two components, so
\[
\dim A
=
\max\{2,1\}
=
\boxed{2}.
\]

This example separates two notions that coincide only under additional equidimensionality
hypotheses: the global dimension of a reducible scheme is the maximum component dimension, not a
number that every irreducible component must share.
\end{solution}

% VI29-DETAILED-CHALLENGE-03
\subsection*{Challenge 3 solution}
\begin{solution}
Let
\[
A=k[x,y,z].
\]
For a prime ideal \(\mathfrak p\), localization at
\[
S=A\setminus\mathfrak p
\]
gives the local ring \(A_{\mathfrak p}\), and VI/29 proves
\[
\dim A_{\mathfrak p}=\operatorname{ht}(\mathfrak p).
\]
Thus it is enough to choose primes of heights \(3,2,1,0\).

Take
\[
\mathfrak p_3=(x,y,z),\qquad
\mathfrak p_2=(x,y),\qquad
\mathfrak p_1=(x),\qquad
\mathfrak p_0=(0).
\]
They are prime because the quotients are respectively
\[
k,\qquad k[z],\qquad k[y,z],\qquad k[x,y,z],
\]
all domains.  The coordinate chain
\[
(0)\subsetneq(x)\subsetneq(x,y)\subsetneq(x,y,z)
\]
shows that their heights are at least \(0,1,2,3\), and the polynomial-ring dimension theorem gives
the matching upper bounds.

For
\[
S_d=A\setminus\mathfrak p_d,
\]
we therefore obtain
\[
\dim S_3^{-1}A=3,\qquad
\dim S_2^{-1}A=2,\qquad
\dim S_1^{-1}A=1,\qquad
\dim S_0^{-1}A=0.
\]
The last localization is
\[
A_{(0)}=\Frac(A)=k(x,y,z).
\]

So a single three-dimensional affine domain contains localizations of every dimension from \(0\)
through \(3\); the dimension retained by localization is controlled by which prime chains survive.
\end{solution}

% VI29-DETAILED-CHALLENGE-04
\subsection*{Challenge 4 solution}
\begin{solution}
Consider the cusp inclusion
\[
A=k[t^2,t^3]\subseteq B=k[t].
\]
VI/28 proved that \(B\) is the normalization of \(A\), so \(B\) is integral over \(A\).

The ring \(B=k[t]\) is a PID of dimension \(1\).  Its prime chains have only the forms
\[
(0)
\]
and
\[
(0)\subsetneq(f),
\]
where \(f\) is irreducible.  Hence their possible lengths are \(0\) and \(1\).

Now contract such a chain to \(A\).  The zero prime contracts to zero because the inclusion
\(A\subseteq B\) is injective.  If \(\mathfrak q=(f)\ne(0)\), its contraction
\[
\mathfrak p=\mathfrak q\cap A
\]
cannot be \((0)\): incomparability for integral extensions would otherwise give comparable primes
\[
(0)\subsetneq\mathfrak q
\]
of \(B\) lying over the same prime \((0)\) of \(A\).  Thus
\[
(0)\subsetneq\mathfrak p
\]
is a strict chain in \(A\).  Since integral extensions satisfy lying over, every prime downstairs
has a prime upstairs, and going up lifts every downstairs chain.

Consequently the attainable prime-chain lengths in \(A\) and \(B\) agree.  Since
\[
\dim B=1,
\]
we get
\[
\boxed{\dim A=1}.
\]
This is the concrete cusp instance of the general theorem
\[
A\subseteq B\text{ integral}
\quad\Longrightarrow\quad
\dim A=\dim B.
\]

At the singular cusp point the local-ring change is
\[
k[t^2,t^3]_{(t^2,t^3)}
\longrightarrow
k[t]_{(t)}.
\]
The rings differ, but the maximal prime-chain length remains \(1\).
\end{solution}

% VI29-DETAILED-CHALLENGE-05
\subsection*{Challenge 5 solution}
\begin{solution}
Let
\[
A=k[x_1,x_2,x_3,\ldots].
\]
For every fixed integer \(n\ge1\), define
\[
\mathfrak p_r=(x_1,\ldots,x_r),
\qquad
1\le r\le n.
\]
Each \(\mathfrak p_r\) is prime because
\[
A/\mathfrak p_r
\cong
k[x_{r+1},x_{r+2},\ldots]
\]
is a domain.  Moreover
\[
(0)\subsetneq
\mathfrak p_1\subsetneq
\mathfrak p_2\subsetneq\cdots\subsetneq
\mathfrak p_n
\]
has length \(n\).

This argument must be read with the quantifiers in the correct order:
\[
\text{for every }n\in\mathbb N,
\quad
\text{there exists a prime chain of length }n.
\]
No one of these finite chains has infinite length.  Rather, the set of attainable lengths contains
\[
\{0,1,2,3,\ldots\},
\]
which is unbounded.  Krull dimension is the supremum of this set, so
\[
\boxed{\dim A=\infty}.
\]

One can also display the infinite ascending sequence
\[
(0)\subsetneq(x_1)\subsetneq(x_1,x_2)\subsetneq\cdots,
\]
but the definition of Krull dimension is already settled by its finite initial segments: the
initial segment ending at \((x_1,\ldots,x_n)\) supplies length \(n\) for every \(n\).
\end{solution}
